\section{Control of long-wavelength excitations}\label{asm:section}\label{sec:assembly}

The local energy estimate controls excitations outside the space of
functions constant on $R$-cells. We now use the two deletion comparisons
to show that the remaining excitations, between that space and the single
constant orbital, have density $o(1)$. The distant comparison makes cell
root-mean-squares nearly constant; the local fourth-moment comparison
relates those root-mean-squares to cell averages. This gives the estimate
needed to pass from local observables to global momentum tests.

\subsection{A variance bound from distant pairs}

We first remove the restriction on the pairs of $R$-cells in
Proposition~\ref{long:comparison}. This is a statement about arrays in
a Hilbert space and has no ground-state assumption.

\begin{lemma}[Distant pairs control the variance]\label{asm:far-pairs}
Let $\mathcal H$ be a complex Hilbert space and let
$f_z\in\mathcal H$, $z\in(\Z/M\Z)^3$. Set
\[
 I_M=\{h\in\{0,\ldots,M-1\}:.30M\le h\le .45M\},
 \qquad \overline f=M^{-3}\sum_z f_z.
\]
For all sufficiently large $M$, there is an absolute constant $C$ such
that
\begin{equation}\label{asm:far-variance}
 \frac1{M^3}\sum_z\|f_z-\overline f\|_{\mathcal H}^2
 \le \frac{C}{M^3|I_M|M^2}
      \sum_z\sum_{h_1\in I_M}\sum_{h_2,h_3\in\Z/M\Z}
             \|f_{z+h}-f_z\|_{\mathcal H}^2.
\end{equation}
\end{lemma}

\begin{proof}
Use the unitary discrete Fourier transform on $(\Z/M\Z)^3$.
The right-hand quadratic form before multiplication by $C$ has Fourier
multiplier
\[
 a_M(k)=\frac1{|I_M|M^2}
   \sum_{h_1\in I_M}\sum_{h_2,h_3}
       |e^{2\pi i k\cdot h/M}-1|^2.
\]
The variance on the left has multiplier one for $k\ne0$ and zero at
$k=0$, so it suffices to bound $a_M(k)$ below away from zero for $k\ne0$.
If $k_2$ or $k_3$ is nonzero modulo $M$, orthogonality gives $a_M(k)=2$.
Otherwise, write $q=\min(k_1,M-k_1)$. For $q\ge1$, the geometric-sum
formula and $\sin(\pi q/M)\ge2q/M$ give
\[
 \left|\frac1{|I_M|}\sum_{h\in I_M}e^{2\pi i k_1h/M}\right|
        \le \frac{C_0}{q}.
\]
Thus $a_M(k)\ge1$ when $q\ge2C_0$. For each of the finitely many
remaining positive $q$, the common value
$a_M(q,0,0)=a_M(M-q,0,0)$ converges to
\[
 \frac1{.15}\int_{.30}^{.45}|e^{2\pi i q t}-1|^2\,\dd t>0.
\]
Their minimum is positive, uniformly for large $M$. Parseval's identity
now proves \eqref{asm:far-variance}, including Hilbert-valued arrays.
\end{proof}

For a fixed nonnegative ground vector $\Phi$ and a cell
$A\in\mathcal C_d\cup\mathcal C_R$, the function $s_A$ of
Lemma~\ref{palm:reference} belongs to
$\mathcal H=L^2(\Lambda_K^{N-1},\dd Y)$. Explicitly,
\begin{equation}\label{asm:rms}
 s_A(Y)=\left(\frac V{|A|}\int_A\Phi(y,Y)^2\,\dd y\right)^{1/2}.
\end{equation}
Applying Lemma~\ref{asm:far-pairs} with $M=K/R$ to
Proposition~\ref{long:comparison} gives
\begin{equation}\label{asm:rms-variance}
 \frac1{|\mathcal C_R|}\sum_{A\in\mathcal C_R}
    \left\|s_A-\frac1{|\mathcal C_R|}\sum_{A'}s_{A'}\right\|_2^2
       =o(D^{-1})+o_K(1).
\end{equation}
The uniformity and order of limits are those of
Section~\ref{set:section}.

\subsection{From root-mean-squares to arithmetic means}

The occupation involves an integral of $\Phi$, whereas $s_A$ is a
root-mean-square of $\Phi$. The next elementary inequality explains why
the fourth-moment local comparison is enough to bridge this difference.

\begin{lemma}[RMS interpolation]\label{asm:rms-interpolation}
For any finite nonempty array of nonnegative numbers $y_i$, with uniform
averaging denoted by $\mathbb E_i$, put
$a=\mathbb E_i y_i^2$ and $m=\mathbb E_i y_i$. Then
\begin{equation}\label{asm:rms-inequality}
 (\sqrt a-m)^2\le\frac12\mathbb E_{i,i'}
           \frac{(y_i-y_{i'})^4}{y_i^2+y_{i'}^2}.
\end{equation}
The quotient is defined to be zero when its denominator is zero.
Furthermore, for every $y,z\ge0$,
\begin{equation}\label{asm:sqrt-comparison}
 |\sqrt y-\sqrt z|^4
 \le\frac{(y-z)^4}{y^2+z^2}
 \le8|\sqrt y-\sqrt z|^4.
\end{equation}
\end{lemma}

\begin{proof}
Let $Q$ denote the double average on the right of
\eqref{asm:rms-inequality}, before its factor $1/2$. Weighted
Cauchy--Schwarz and independence of the two uniform indices give
\[
 [2(a-m^2)]^2
 =\bigl[\mathbb E_{i,i'}(y_i-y_{i'})^2\bigr]^2
 \le Q\mathbb E_{i,i'}(y_i^2+y_{i'}^2)=2aQ.
\]
If $a>0$, divide by $4(\sqrt a+m)^2\ge4a$; the case $a=0$ is
immediate. For the second statement factor
$y-z=(\sqrt y-\sqrt z)(\sqrt y+\sqrt z)$. The remaining ratio is
$(\sqrt y+\sqrt z)^4/(y^2+z^2)$, which lies between $1$ and $8$.
\end{proof}

For any two $d$-cells $z,z'$ write
\[
 Q_{z,z'}(Y)=
   \frac{|s_z(Y)-s_{z'}(Y)|^4}{s_z(Y)^2+s_{z'}(Y)^2}.
\]
For a nearest-neighbor edge $e=(z,z')$, abbreviate this as $Q_e$.
Inside one $R$-cell, any two $d$-cells can be joined by a path of at most
$3R/d$ nearest-neighbor steps that stays in the cell. Put
$t_z(Y)=\sqrt{s_z(Y)}$. If the path has $m$ edges, then
\eqref{asm:sqrt-comparison} and H\"older's inequality give
\[
 Q_{z,z'}\le8|t_z-t_{z'}|^4
       \le8m^3\sum_{e=(v,w)\text{ on the path}}|t_v-t_w|^4
       \le8m^3\sum_{e\text{ on the path}}Q_e.
\]
Enlarging the last sum to all internal edges yields
\begin{equation}\label{asm:path-fourth}
 Q_{z,z'}(Y)\le C_jR^3
       \sum_{e\text{ internal edge of the }R\text{-cell}}Q_e(Y).
\end{equation}
The full internal edge sum bounds every pair, so no estimate of path
multiplicities is needed. Averaging over pairs and over $R$-cells, and using
$|\mathcal C_d|/|\mathcal C_R|=(R/d)^3$, consequently gives
\begin{align}
 &\frac1{|\mathcal C_R|}\sum_{A\in\mathcal C_R}
   \frac1{|\{z\in\mathcal C_d:z\subset A\}|^2}
    \sum_{z,z'\subset A}\int Q_{z,z'}\,\dd Y
 \notag\\
 &\hspace{20mm}\le C_jR^6\,
       \frac1{|\mathcal C_d|}\sum_{e\text{ a nearest-neighbor edge}}
             \int Q_e\,\dd Y
 \notag\\
 &\hspace{20mm}\le C_j R^6 D^{-1-u/8}+o_K(1)
 =o(D^{-1})+o_K(1).
 \label{asm:all-local-pairs}
\end{align}
The second inequality is Proposition~\ref{palm:local-comparison}.
The last uses $6w<u/8$, which follows from \eqref{set:exponents}.
For fixed dilution the factor $R^6$ is fixed, so it causes no difficulty
for the thermodynamic remainder.

Define
\begin{equation}\label{asm:field-means}
 F(y,Y)=\sqrt V\Phi(y,Y),\qquad
 m_A(Y)=\frac1{|A|}\int_A F(y,Y)\,\dd y.
\end{equation}
Spatial integrals of $F$ over the whole torus will use the normalized
measure $\dd y/V$. For each cell, nonnegativity of $F$ implies
$0\le m_A\le s_A$, and
\[
 (s_A-m_A)^2\le s_A^2-m_A^2
       =\frac1{|A|}\int_A|F-m_A|^2\,\dd y.
\]
It follows from \eqref{en:d-occupation} and bosonic symmetry that
\begin{equation}\label{asm:micro-mean-rms}
 \frac1{|\mathcal C_d|}\sum_{z\in\mathcal C_d}
       \|m_z-s_z\|_2^2
 \le\|(W_d)_1\Phi\|_2^2
 =\frac1D\frac1V\langle\mathrm d\Gamma(W_d)\rangle_\Phi
 \le D^{-1}(\epsilon_j+o(1))+o_K(1),
\end{equation}
where $\epsilon_j\to0$ as $j\to\infty$, and the dilute remainder is taken
at fixed $j$.

For an $R$-cell $A$, the number $s_A(Y)$ is the root-mean-square of
the numbers $s_z(Y)$ over its $d$-cells. Let $\overline s_A(Y)$ denote
their arithmetic mean. Lemma~\ref{asm:rms-interpolation} and
\eqref{asm:all-local-pairs} imply
\begin{equation}\label{asm:R-rms-arithmetic}
 \frac1{|\mathcal C_R|}\sum_A\|s_A-\overline s_A\|_2^2
           =o(D^{-1})+o_K(1).
\end{equation}
Also $m_A$ is the arithmetic mean of the $m_z$ in $A$. Jensen's
inequality, \eqref{asm:micro-mean-rms}, and
\eqref{asm:R-rms-arithmetic} therefore give
\begin{equation}\label{asm:R-mean-rms}
 \frac1{|\mathcal C_R|}\sum_A\|m_A-s_A\|_2^2
       \le D^{-1}(2\epsilon_j+o(1))+o_K(1).
\end{equation}

\subsection{The infrared estimate and arbitrary ground vectors}

\begin{proposition}[Remaining excitations]\label{asm:infrared}
Put $Q_R=P_R-P_0$. Uniformly over normalized, possibly complex,
bosonic ground vectors,
\begin{equation}\label{asm:infrared-bound}
 D\|(Q_R)_1\Phi\|_2^2\le16\epsilon_j+o(1)+o_K(1),
\end{equation}
where $\epsilon_j\to0$. The coefficient $16$ is independent of $j$;
the dilute remainder may depend on $j$. The exponents $u,w$ and the scale $R$ may be
kept fixed when $j$ is increased; all dilute and volume thresholds may
depend on $j$.
\end{proposition}

\begin{proof}
First let $\Phi\ge0$. For Hilbert-valued arrays, the variance is the
minimum of the mean squared distance to a constant array. Comparing the
array $m_A$ with the mean of the array $s_A$ thus yields
\[
 \frac1{|\mathcal C_R|}\sum_A\|m_A-\overline m\|_2^2
 \le\frac2{|\mathcal C_R|}\sum_A\|m_A-s_A\|_2^2
      +\frac2{|\mathcal C_R|}\sum_A\|s_A-\overline s\|_2^2.
\]
Here bars without cell subscripts denote averages over all $R$-cells.
The left side is exactly $\|(P_R-P_0)_1\Phi\|_2^2$, by
\eqref{asm:field-means} and the normalized spatial measure.
Equations~\eqref{asm:rms-variance} and \eqref{asm:R-mean-rms} prove
\eqref{asm:infrared-bound} for nonnegative vectors, with coefficient $4$
in place of $16$. The multiplication
of a volume remainder by $D$ is harmless in the inner limit because
$D\to D_*$.

For a real bosonic ground vector $f$, its positive and negative parts
$f_+,f_-$ belong to the Dirichlet form domain and remain symmetric.
The Sobolev truncation identities give
\[
 \|f\|_2^2=\|f_+\|_2^2+\|f_-\|_2^2,
 \qquad q[f]=q[f_+]+q[f_-].
\]
Each nonzero part has Rayleigh quotient at least $E_N$, and their sum
attains $E_N$. Both nonzero parts therefore attain $E_N$ and are ground
vectors. Complex conjugation commutes with the operator, so the real
and imaginary parts of a complex ground vector are ground vectors or
zero. A normalized $\Phi$ can consequently be written as
\[
 \Phi=\sum_{a=1}^m c_a\Phi_a,
 \qquad m\le4,\quad \Phi_a\ge0,\quad\|\Phi_a\|_2=1,
 \qquad \sum_{a=1}^m|c_a|^2=1,
\]
where each $\Phi_a$ is a ground vector. Cauchy--Schwarz gives
\[
 \|(Q_R)_1\Phi\|_2^2
       \le\sum_{a=1}^m\|(Q_R)_1\Phi_a\|_2^2.
\]
Thus at most a factor four is lost in the vanishing error in
\eqref{asm:infrared-bound}. No comparison of component phases or
connectedness assumption is used.
\end{proof}

Because $P_0\le P_R$, $Q_R$ and $W_R$ are orthogonal projections with
\begin{equation}\label{asm:orthogonal-decomposition}
 1-P_0=W_R+Q_R.
\end{equation}
Consequently every normalized bosonic ground vector satisfies the exact
identity
\begin{equation}\label{asm:depletion-decomposition}
 D(1-B(\Phi))=
        \frac1V\langle\mathrm d\Gamma(W_R)\rangle_\Phi
                   +D\|(Q_R)_1\Phi\|_2^2.
\end{equation}
We record the conclusions in the per-volume normalization used for
momentum observables:
\begin{align}
 \frac1V\langle\mathrm d\Gamma(P_R-P_0)\rangle_\Phi
   &=D\|(P_R-P_0)_1\Phi\|_2^2=o(1),
       \label{asm:coarse-density}\\
 \frac1V\langle\mathrm d\Gamma(1-P_0)\rangle_\Phi
   &=D(1-B(\Phi))=\cD\alpha^{3/2}+o(1).
       \label{asm:scaled-depletion}
\end{align}
Both remainders are uniform over normalized, possibly complex ground
vectors in the iterated limit of Section~\ref{set:section}. For the first,
bosonic symmetry gives the equality, and Proposition~\ref{asm:infrared}
gives the error estimate: choose $j$ large and then increase the lower
bound on $D$ and the volume threshold. The left side is independent of
$j$, since $u,w,R$ remain fixed. The second follows from
\eqref{asm:depletion-decomposition} and \eqref{en:R-occupation}, whose
energy estimate applies directly to complex ground vectors.
